\chapter{Hip\'otesis}

El desgaste de los motores de aeronaves y el deterioro de su rendimiento durante la fase de despegue puede evaluarse y detectarse de forma precisa sin desmontar el motor ni requerir conocimientos expertos, mediante la implementaci\'on de modelos de aprendizaje autom\'atico (ML).

En particular, se plantea que, usando algoritmos espec\'ificos de ML como la Regresi\'on de Procesos Gaussianos (GPR), las M\'aquinas de Vectores de Soporte (SVM) y los Perceptrones Multicapa (MLP), entrenados con datos operativos y ambientales de la base de datos de emisiones de motores de la OACI (EEDB), es posible predecir el par\'ametro de flujo de combustible en despegue (\textit{Fuel Flow T/O}). A partir de esta predicci\'on, se puede establecer un intervalo de confianza estad\'istico (al 99 \%) que delimite los rangos de operaci\'on normal; de modo que, si el valor medido se encuentra fuera de este intervalo, se identifique de forma autom\'atica una anomal\'ia o deterioro en el rendimiento del motor.
